\documentclass[25pt]{tikzposter} % See Section 3
%25pt,margin=0mm,innermargin=15mm
\usepackage{amsmath,graphicx}
\newcommand{\hsc}[1]{{\LARGE\MakeUppercase{#1}}}
\newcommand{\hsct}[1]{{\Large\MakeUppercase{#1}}}

%\makeatletter
%\input{theguy28pt.clo}
%\makeatother

\title{\bf Digital Predistortion with Low Precision ADCs} 
\institute{Rice University} % See Section 4.1
%\titlegraphic{\includegraphics[width=5in]{Rice_logo_hi_res.png}}
\author{Chance Tarver and Joseph R. Cavallaro} 
\usetheme{Basic} % See Section 5

%% THEME STUFF %%
\geometry{paperwidth=42in,paperheight=42in}
\definecolor{riceblue}{RGB}{0,36,106}
\definecolor{ricegray}{RGB}{94,96,98}
%\titlegraphicheight{1mm}
\tikzposterlatexaffectionproofoff  %Turn off watermark thing in corner

%Put logo at the side
\definetitlestyle{sampletitle}{
width=40in, roundedcorners=20, linewidth=2pt, innersep=5pt,
titletotopverticalspace=15mm, titletoblockverticalspace=18mm,
titlegraphictotitledistance = 0pt
}{
\begin{scope}[line width=\titlelinewidth, rounded corners=\titleroundedcorners]
\draw[color=blocktitlebgcolor, fill=titlebgcolor]
(\titleposleft,\titleposbottom) rectangle (\titleposright,\titlepostop);
\end{scope}
}

 \usetitlestyle{sampletitle}
%Recalculate the size of the textwidth
\makeatletter
\setlength{\TP@visibletextwidth}{\textwidth-2\TP@innermargin}
\setlength{\TP@visibletextheight}{\textheight-2\TP@innermargin}
\makeatother

%Set the colors
\colorlet{backgroundcolor}{riceblue}
\colorlet{framecolor}{black}
\colorlet{blocktitlebgcolor}{ricegray} %<---- change color
\colorlet{titlebgcolor}{white}
\colorlet{titlefgcolor}{black}  %Title font color
\colorlet{notebgcolor}{ricegray}
\colorlet{notefgcolor}{white}
\colorlet{noteframecolor}{black}

\makeatletter
\usepackage{helvet}
\renewcommand*{\familydefault}{\sfdefault}% Let's have a sans serif font - See more at: http://latex-cookbook.net/cookbook/examples/poster/#sthash.jDLWH62B.dpuf
\renewcommand\TP@maketitle{%
  \centering
   \begin{minipage}[b]{1\linewidth}
        \centering
        \vspace*{1em}
        \color{black} %color of letters
        {\bfseries \Huge \sc \@title \par}
        \vspace*{1em}
        {\huge \@author \par}
        \vspace*{1em}
        {\LARGE \@institute}
    \end{minipage}%
     \tikz[overlay]\node[scale=1,anchor=east,xshift=16in,yshift=1.5in,inner sep=0pt] {%
       \includegraphics[width=6in]{ECE}
          }; 
      \tikz[overlay]\node[scale=1,anchor=east,xshift=4.4in,yshift=1.6in,inner sep=0pt] {%
       \includegraphics[width=6in]{Rice_logo_hi_res.png}       
   };   
}
\makeatother
\makeatletter
\input{theguy30pt.clo}
\makeatother

\begin{document}


\maketitle % See Section 4.1

\begin{columns} % See Section 4.4
\centering
\column{0.33} % See Section 4.4
\block{Motivation}
{    \Large
\begin{itemize}
     \item {\bf Spectrum Scarcity $\rightarrow$ Frequency Agile Standards}
            \begin{itemize}
                \item Non-contiguous Transmission
                \item Carrier Aggregation (CA) in LTE-Advanced
                \item Cognitive Radio
                \item 5G Cellular
            \end{itemize}
        \item {\bf Non-contiguous carriers intermodulate}
        \begin{itemize}
            \item Caused by nonlinearities in power amplifiers (PAs)
            \item Undesired spurious emissions (spurs)
            \item Could interfere with nearby channels
            \item Self-interference to own receiver when using FDD
            %frequency-division duplexing
        \end{itemize}
        \item{\bf Current 4G chipsets support up to 4 carriers}
            \begin{itemize}
                \item Snapdragon 835
                \item 4x20 MHz carrier aggregation downlink, 2x20 MHz uplink
            \end{itemize}
            \vspace{1mm}
        \item{\textbf{\textit{ Need efficient way to linearize for this scenario}}}
        \item{\textbf{\textit{ DPD requires extra hardware}}}
        	\begin{itemize}
        		\item Extra RX chains
        		\item Larger area 
        		\item More Power
        	\end{itemize}
     \end{itemize}  }
     \block{Power Spectral Density}{
%\begin{tikzfigure}[\large Intermodulations when broadcasting more than 2 carriers.]
%\label{fig:fig1}
%\includegraphics[width=0.95\linewidth]{SubBandIM3s2}
%\end{tikzfigure}
}   
\block{Related Works}{
\Large
    \begin{itemize}
        \item {\bf Reduce Power}
            \begin{itemize}
                \item Operate in a more linear PA region
                \item Less range and less power efficient
            \end{itemize} 
                        \vspace{1mm}
        \item {\bf Full-Band Digital Predistortion (DPD)}
            \begin{itemize}
                \item Computationally expensive
                \item Does not scale for noncontiguous carriers
                \item Requires large sampling rate as carrier spacing grows
            \end{itemize}  
                        \vspace{1mm}
        \item {\bf Sub-Band DPD}
            \begin{itemize}
                \item Previously explored by the authors with the WARP SDR RF Board
                \item Observes and applies DPD to individual spurs
                \item Can reduce the necessary sampling rate and complexity
                \item Has only been considered for 2 carriers           
            \end{itemize}                 
    \end{itemize}
}
\block{Main Idea}{\Large
\begin{itemize}
    \item {\bf Learn DPD coefficient, $\alpha$, for each spur}
        \begin{itemize}
            \item Iteratively learn coefficients as necessary using adaptive, LMS algorithm.
            \item Apply them as in Equation 5 to reduce spurious emissions.
        \end{itemize}
\end{itemize}
}        
     
     

\column{0.33}
\block{Mathematical Model}{\Large
\begin{itemize}
    \item {\bf PA Inputs and Outputs:}
    \begin{itemize}
        \item Parallel Hamerstein baseband PA model
        \item Third order, memoryless
    \end{itemize}
\begin{align}
\text{\bf Input:} \ x(n) =& \sum_{i = 1}^N x_i(n) e^{2\pi j n \frac{f_i}{f_s}}, \label{eq:PA_In} \\
\text{\bf Output:} \ y(n) =& \beta_1 x(n) + \beta_3 |x(n)|^2 x(n). \label{eq:PA}
\end{align}
\vspace{1mm}
\item {\bf IM3 Spurious Signals and Their Locations}
\begin{align}
y_{IM3:i;j,k}(n) &= x_i^*(n)x_j(n)x_k(n), \label{eq:IM3_out} \\
f_{IM3:i;j,k} &= -f_i + f_j + f_k.
\end{align}
\item {\bf DPD Processing}
\begin{align}
\tilde{x}(n) &= 
\sum_{i=1}^N x_i(n) e^{2\pi j n \frac{f_i}{f_s} n}  \nonumber\\ 
&+ \sum_{i = 1}^{N} \sum_{j = 1}^N \sum_{k = 1}^N \alpha_{i;j,k}  \ x_i^*(n) x_j(n) x_k(n) e^{2\pi j n\frac{-f_i + f_j + f_k}{f_s}}.
\label{eq:PA_In_dpd}
\end{align}
\end{itemize}
}


\block{M\hsct{atlab} Simulation}{\Large
\begin{itemize}
    \item {\bf LTE-Advanced CA Scenario}
    \begin{itemize}
        \item Three, 20MHz CCs, scenario CA 41C-41A
        \item Bandpass filter around passpand on the RF Frontend
        \item Only one spurious region is in violation of the emission limits 
    \end{itemize}
    \vspace{1mm}
    \item {\bf Result}
    \begin{itemize}
        \item Rapid convergence of coefficients
        \item Suppression below the emission limit in the spurious region
    \end{itemize}
\end{itemize}
%\begin{tikzfigure}[]
%\label{fig:fig1}
%\centering
%\includegraphics[]{filename2}
%\end{tikzfigure}

%\begin{tikzfigure}[]
%\label{fig:fig1}
%\includegraphics[]{filename}
%\end{tikzfigure}
}


\column{0.33}
\block{DPD Coefficient Learning}{\Large
%\begin{tikzfigure}
%\includegraphics[width=0.97\linewidth]{SystemArchi}
%\end{tikzfigure}
\vspace{3mm}
\begin{itemize}
    \item {\bf LMS Adaptive Training}
        \vspace{2mm}
           \begin{align} 
%    \boldsymbol{\alpha}(m) &= [\alpha_0(m) \:\: \alpha_2(m) \:\: ... \:\: \alpha_N(m)]^T \\
    \alpha_{i;j,k}(n+1) &= \alpha_{i;j,k}(n) -\mu \frac{ x_i^*(n) x_j(n) x_k(n){e}^*_{i;j,k}(n)}{|| x_i^*(n) x_j(n) x_k(n)||}
%    \label{eq:NLMS_Update_memory}
    \end{align}
    \vspace{1mm}
    \begin{itemize}
        \item Decorrelates the error signal, $e$, with basis functions from Equation 3 to learn the DPD coefficients, $\alpha$.
        \item Rate is controlled by LMS parameter $\mu$.
        \item Learning is resilient to temperature changes, PA Gain changes, and changes in signal characteristics.
    \end{itemize}
\end{itemize}

}

\block{G\hsct{nu}Radio Simulator}{\Large
%\begin{itemize}
    %\item {\bf GNURadio}
    %\begin{itemize}
        %\item Software defined radio (SDR) development platform
        %\item Can work with SDR boards such as the USRP
        %\item Efficient, parallel, real-time CPU implementation 
        %\item Uses python and C++ for easy development and performance 
    %\end{itemize}
    %\vspace{1mm}
%\item{\bf DPD Simulator}
    %\begin{itemize}
        %\item Custom GNURadio flowgraph, modules, and blocks for performing DPD
        %\item Movable component carriers
        %\item Changeable PA model
        %\item Can add and remove DPD processing
    %\end{itemize}
%    \begin{tikzfigure}
%\includegraphics[width=0.98\linewidth]{gnu.png}
%\end{tikzfigure}

%\end{itemize}
}
\note[targetoffsetx=0.5in,targetoffsety=-5.5in,width = 6in,connection,rotate=2]{\bf Can change carrier magnitude and placment, PA model, and DPD application dynamically.}



\block{Future Work}{\Large
\begin{itemize}
    \item {\bf Main carrier linearization}
    %\begin{itemize}
    %    \item Use a variation of the sub-band DPD method to reduce in-band emission. 
    %\end{itemize}
    \item {\bf Hardware testing on the USRP with external PAs }
    %\begin{itemize}
    %    \item USRP SDR RF Board
    %    \item Off-the-shelf, external PAs
    %\end{itemize}
\end{itemize}
}

\end{columns}
\end{document}